在平方根归一化的三歧相变之外，RH-sharp 临界点本身还允许进一步压缩出一个单步代数滤子与一个余维一的误差指数尖点。

\begin{theorem}[临界平方根模的一阶湮灭]\label{thm:xi-time-part65-critical-sqrt-mode-first-order-annihilation}
\leanverified{paper\_xi\_time\_part65\_critical\_sqrt\_mode\_first\_order\_annihilation}
沿用命题 \ref{prop:real-input-40-collision-rhk-window} 的记号，设
\[
u_\ast:=u_R,
\qquad
s:=-\lambda_-(u_\ast)>0.
\]
则
\[
\lambda_+(u_\ast)=s^2,
\qquad
\lambda_-(u_\ast)=-s.
\]
进一步，对任意形如
\[
a_n=A\,\lambda_+(u_\ast)^n+B\,\lambda_-(u_\ast)^n+r_n
\qquad (n\ge 0)
\]
的标量序列，只要存在 \(\eta<s\) 使
\[
r_n=O(\eta^n),
\]
则一阶滤子
\[
\mathcal L_s a_n:=a_{n+1}+s\,a_n
\]
满足
\[
\mathcal L_s a_n
=(s^2+s)A\,s^{2n}+\widetilde r_n,
\qquad
\widetilde r_n=O(\eta^n)=o(s^n).
\]
特别地，负谱模 \(\lambda_-(u_\ast)^n\) 被 \(\mathcal L_s\) 单步精确湮灭，而保留下来的主导项只来自临界 Perron 模 \(s^{2n}\)。
\end{theorem}
\begin{proof}
注 \ref{rem:collision-cubic-rhsharp-quintic} 已给出
\[
\lambda_+(u_\ast)=s^2,
\qquad
\lambda_-(u_\ast)=-s.
\]
因此
\[
\mathcal L_s\bigl(\lambda_-(u_\ast)^n\bigr)
=(-s)^{n+1}+s(-s)^n
=0,
\]
而
\[
\mathcal L_s\bigl(\lambda_+(u_\ast)^n\bigr)
=s^{2n+2}+s^{2n+1}
=(s^2+s)s^{2n}.
\]
若 \(r_n=O(\eta^n)\) 且 \(\eta<s\)，则
\[
\mathcal L_s r_n=r_{n+1}+s\,r_n=O(\eta^n)=o(s^n).
\]
将上述三式按线性叠加，便得到所示展开。证毕。
\end{proof}

\begin{corollary}[平方根律的尖点接触与非厚相性]\label{cor:xi-time-part65-sqrt-law-cusp-nonthick-phase}
沿用推论 \ref{cor:real-input-40-collision-error-phase} 的记号，定义
\[
\beta(u):=\frac{\log \Lambda(u)}{\log \lambda_+(u)}.
\]
则
\[
\beta(u)=0
\qquad (0<u\le 1),
\]
\[
\beta(u)\le \frac12
\qquad (1<u\le u_R),
\]
且等号仅在 \(u=u_R\) 发生；此外，
\[
\beta(u)>\frac12
\qquad (u>u_R),
\qquad
\beta(u)\to 1
\qquad (u\to\infty).
\]
因此，\(\beta=\frac12\) 的平方根级误差并不形成参数上的开相区，而只在唯一临界参数 \(u_R\) 处与相图发生余维一接触。
\end{corollary}
\begin{proof}
推论 \ref{cor:real-input-40-collision-error-phase} 已逐段给出 \(\beta(u)\) 在
\((0,1]\)、\((1,u_R]\) 与 \((u_R,\infty)\) 上的精确刻画，并证明
\(\beta(u)\to 1\) 当 \(u\to\infty\)。将其直接合并，便得到全部结论。证毕。
\end{proof}

\endinput
